\section{General principles for graded lattice separation}
Let $Q_1,Q_2$ be positive-definite real quadratic forms on $\mathbb Z^d$, with an integer grading $g$. Let $N_i(e,k)$ count charges at energy $e$ and grade $k$, and put
\[
D(Q_1,Q_2;g)=\inf\{e\ge0:\exists k,\ N_1(e,k)\ne N_2(e,k)\},
\]
with $\inf\varnothing=+\infty$.

\begin{proposition}[Attainment and relabelling invariance]
If $D(Q_1,Q_2;g)<\infty$, its infimum is attained. If a bijection $f:\mathbb Z^d\to\mathbb Z^d$ obeys $Q_2(f(v))=Q_1(v)$ and $g(f(v))=g(v)$, then $D(Q_1,Q_2;g)=+\infty$.
\end{proposition}
\begin{proof}
For $i=1,2$, positive definiteness gives $Q_i(v)\ge c_i\|v\|^2$ for $c_i>0$. Every bounded energy interval therefore contains finitely many charge states from either system. A nonempty set of differing energies has a least element. The claimed bijection identifies each graded energy-level set and therefore every multiplicity.
\end{proof}

\begin{proposition}[Conditional local witness persistence]
Let $Q_{\sigma,x}$, $\sigma=1,2$, be continuous positive-definite quadratic-form families near $x_0$, uniformly coercive as $Q_{\sigma,x}(v)\ge c\|v\|^2$ for fixed $c>0$. Assume that finite distinguished witness sets have a common continuous energy $T(x)$ near $x_0$, every other nonzero-grade charge has energy strictly greater than $T(x_0)$ at $x_0$, the refined counts at $T(x)$ differ throughout a neighborhood, and the grade-zero multiplicities of the two systems agree at all energies throughout that neighborhood. Then $D(Q_{1,x},Q_{2,x};g)=T(x)$ on a possibly smaller neighborhood.
\end{proposition}
\begin{proof}
Shrink to a compact parameter neighborhood where $T$ is bounded. Uniform coercivity confines all charges with energy at most $\sup T+1$ to a finite lattice ball. By continuity, all finitely many strict competitor inequalities persist on a smaller neighborhood. Thus no nonzero-grade charge occurs below $T(x)$. The assumed grade-zero matching excludes lower differences, and the witness imbalance at $T(x)$ supplies the first difference.
\end{proof}
